\section{Validation: from RMSE to decision-aware metrics}
\label{sec:validation}

Most surrogate-modelling studies in energy-system optimization validate
their models with test-set RMSE, MAE and $R^{2}$: they check how
accurately the surrogate reproduces known outputs on held-out data.
If those scores are acceptable, the trained model is typically placed
inside an optimizer, without further tests on the quality of the
decisions it produces there.

We argue that this workflow is a sensible first step, but not enough
on its own. A surrogate can predict well on average and still yield
poor optimization outcomes---plans that violate constraints, are
clearly suboptimal compared with the full model, or change sharply
when inputs shift slightly. RMSE-type metrics average prediction
errors uniformly, whereas an optimizer reacts to errors in the
constraints that actually bind and in regions where, for probabilistic
surrogates, the posterior---the predicted mean and spread at each
input---is locally biased.

The subsections below review validation practice under six headings:
standard point and interval metrics (Section~\ref{sec:val:point});
feasibility and violation severity (Section~\ref{sec:val:feasibility});
decision-aware cost and gap metrics (Section~\ref{sec:val:decision});
stress tests under extreme operating conditions
(Section~\ref{sec:val:stress}); out-of-distribution and policy-shift
evaluation (Section~\ref{sec:val:ood}); and reproducibility of
reported results (Section~\ref{sec:val:repro}). The corresponding
metrics and test designs are summarized in
Table~\ref{tab:T4-validation}.

% BEGIN inlined table table_T4_validation
\begin{table*}[!t]
  \centering
  \footnotesize
  \caption{Validation metrics used to qualify surrogates in
  energy-system optimization. The first block lists the standard
  regression metrics; the second block lists decision-aware metrics
  that we recommend reporting alongside them; the third block lists
  the test designs that distinguish a healthy validation from a
  cosmetic one. References are representative; the full study-level
  picture is in Table~\ref{tab:T6-evidence-map}.}
  \label{tab:T4-validation}
  \renewcommand{\arraystretch}{1.2}
  \begin{tabularx}{\textwidth}{@{}L{0.22\textwidth} L{0.16\textwidth} L{0.26\textwidth} L{0.28\textwidth}@{}}
    \toprule
    Metric / test & Tells you about & When it is useful (refs.) & Common misuses (refs.) \\
    \midrule
    \multicolumn{4}{l}{\textit{Standard regression metrics}} \\
    RMSE / MAE / $R^2$
      & Average fit
      & Sanity / first screen
        \cite{Mylonopoulos202332697,Ruan2021221,Starke2025214,Perera2019191}
      & Reported alone, without decision-aware companion
        \cite{Khaloie2025,Lim2025,Mylonopoulos202332697} \\
    Quantile / prediction-interval coverage
      & Calibration of posterior spread
      & Probabilistic surrogates, BO acquisition
        \cite{Hu2020,Hu2021671,Volodina2022,Hu2024,Buchanan2024}
      & Optimistic on training set; calibration breakdown OOD
        \cite{Subedi2026157,Khaloie2025} \\
    \midrule
    \multicolumn{4}{l}{\textit{Decision-aware metrics}} \\
    Decision regret (cost gap to oracle)
      & Loss in decision quality
      & Surrogate inside optimizer (P1, P3)
        \cite{Chen20244723,Wu2022231,Liang20246508,Wu2024391}
      & Computed only on training set / on synthetic data
        \cite{Ruan2021221,Khaloie2025,Lim2025} \\
    Feasibility rate
      & Constraint safety on the host problem
      & Hard-constrained problems (P1, P3, P5)
        \cite{Liu20242534,Chen2023657,Chen2022,Koirala20242284}
      & Reported without violation severity
        \cite{Khaloie2025,Lim2025,Starke2025214} \\
    Violation severity (worst case + 95th percentile)
      & Robustness under tail load / contingencies
      & Safety-critical operation, stress regimes
        \cite{Chen2023657,Liang20246508,Chen20244723}
      & Average-only reporting; missing contingency stratification
        \cite{Lim2025,Khaloie2025} \\
    Objective gap distribution (MIP gap, CVaR of gap)
      & Distributional fairness across instances
      & MILP / MINLP surrogate use
        \cite{Wu2022231,Chen20244723,Wu2024391}
      & Mean-only summaries hide tail behaviour
        \cite{Ruan2021221,Mylonopoulos202332697} \\
    Out-of-distribution score
      & Trust-region applicability
      & Online updating, planning under future capacity mixes
        \cite{Park2019,Subedi2026157,Liu20222679}
      & Single fixed test split, no policy-shift block
        \cite{Khaloie2025,Lim2025,Starke2025214} \\
    \midrule
    \multicolumn{4}{l}{\textit{Test designs}} \\
    Stress tests (peak load, low-renewable spell, cold spell)
      & Tail behaviour
      & ESM operation under extremes
        \cite{Hu2024,Du2025,Yang2024__2,Bayani20231024}
      & Only on average days; no Dunkelflaute / heat-wave block
        \cite{Khaloie2025,Lim2025,Starke2025214} \\
    Counterfactual policy shift (capacity mix, tariff regime)
      & Robustness to scenario change
      & Capacity planning, MES design under scenarios
        \cite{Jalving2023,Park2019}
      & Validation only on historical mix
        \cite{Khaloie2025,Lim2025,Mylonopoulos202332697} \\
    Scenario-shift / distribution-shift split
      & Generalisation beyond IID
      & Stochastic / robust use, OOD planning
        \cite{Volodina2022,Lawson201647,Subedi2026157,Liu20222679}
      & IID random split only
        \cite{Khaloie2025,Lim2025,Starke2025214} \\
    Holdout reproducibility (seeds, splits, code, data)
      & Trust and comparability
      & All surrogate classes
        \cite{Khazeiynasab2021,Zhang2025__7}
      & Single seed / split; missing code or data release
        \cite{Khaloie2025,Lim2025,Starke2025214,Ruan2021221,Mylonopoulos202332697} \\
    \bottomrule
  \end{tabularx}
\end{table*}
% END inlined table table_T4_validation

\subsection{Point and interval prediction metrics}
\label{sec:val:point}

Point and interval metrics remain the entry-level validation step.
Given a held-out test set of $n$ pairs $(\mathbf{x}_i, y_i)$, where
$y_i$ denotes a quantity returned by the high-fidelity model and
$\hat{y}_i$ the surrogate prediction, the most common point metrics are
\begin{align}
  \mathrm{MAE}  &= \frac{1}{n}\sum_{i=1}^{n}\lvert y_i - \hat{y}_i\rvert,
  \label{eq:mae} \\
  \mathrm{RMSE} &= \sqrt{\frac{1}{n}\sum_{i=1}^{n}(y_i - \hat{y}_i)^2},
  \label{eq:rmse} \\
  R^{2}         &= 1 - \frac{\sum_{i=1}^{n}(y_i - \hat{y}_i)^2}
                       {\sum_{i=1}^{n}(y_i - \bar{y})^2},
  \label{eq:r2}
\end{align}
with $\bar{y} = \frac{1}{n}\sum_i y_i$.
MAE averages absolute errors and is easy to interpret in the units of
$y$; RMSE penalizes large errors more heavily; $R^{2}$ measures how
much variance in $y$ is explained relative to a constant baseline.
These scores are reported in most surrogate studies covered by this
review; they are the natural first screen when a surrogate supplies
forecasts or metamodel outputs to a downstream planner
\cite{Ruan2021221,Khaloie2025,Lim2025}.
Because the metric set itself is not discriminating, we cite only
representative energy-system cases that differ in domain or surrogate
class: polynomial response surfaces for integrated electric--thermal
scheduling ($R^{2}$; \cite{Wang2020202933}), fused ML models for
regional integrated energy system loads (RMSE;
\cite{Shi2024__2}), LSTM surrogates for district heating with thermal
storage (MAE/RMSE; \cite{Du2025}), Gaussian-process surrogates in
heat--electricity microgrids (MAE; \cite{Li2023__3}), and neural-
network surrogates for distributed energy-system and multi-energy-
storage optimization (MAE/$R^{2}$; \cite{Perera2019191,Ren2024}).

For probabilistic surrogates, validation additionally requires that the
reported uncertainty band matches empirical coverage.
A common check is the prediction-interval coverage probability
(PICP): the fraction of test points that fall inside a
$(1-\alpha)$ prediction interval $[\hat{y}_i^{\mathrm{lo}},
\hat{y}_i^{\mathrm{hi}}]$,
\begin{equation}
  \mathrm{PICP} = \frac{1}{n}\sum_{i=1}^{n}
  \mathbb{I}\!\left[\hat{y}_i^{\mathrm{lo}} \le y_i \le
  \hat{y}_i^{\mathrm{hi}}\right],
  \label{eq:picp}
\end{equation}
where $\mathbb{I}[\cdot]$ is the indicator function.
Pinball loss and the continuous ranked probability score (CRPS)
summarize quantile and full predictive-distribution accuracy,
respectively \cite{Hu2024,Sun20191497}.
Representative interval checks include Transformer--GP net-load
forecasts (PICP, CRPS; \cite{Hu2024}), GPR renewable intervals for
microgrid robust scheduling (coverage probability; \cite{Yang2024__2}),
linked GP emulators in municipal energy planning
\cite{Volodina2022,Lawson201647}, and GP emulation of stochastic
dispatch under wind uncertainty \cite{Hu2020,Hu2021671}.

These metrics are necessary but not sufficient.
CNN-based inverter surrogates can retain low bulk RMSE while interval
calibration breaks down in operating regions absent from the training
set \cite{Subedi2026157}.
The literature therefore treats point and interval scores as a screening
filter rather than as a final acceptance criterion for surrogate-
assisted optimization \cite{Khaloie2025,Lim2025,Starke2025214}.

\subsection{Constraint feasibility and violation severity}
\label{sec:val:feasibility}

When the surrogate is used inside the optimizer (P1, P3, P5),
feasibility matters more than the point and interval prediction metrics
described above because the key question is not how well it predicts on
average, but whether the operating plan or design it produces is
actually feasible. The primary validation metric is therefore the
feasibility rate $\phi_{\mathrm{feas}}$ (Eq.~\eqref{eq:feas-rate}).
Here a \emph{test scenario} is one operating condition of the
optimization problem that was excluded from surrogate training---for
example one day of load and renewable profiles, one network loading
case, or one draw from an uncertainty set. For each of $N$ such
scenarios, the surrogate-based optimizer is run and returns an
\emph{operating decision} $\mathbf{x}_j$: the dispatch, unit
commitment, storage schedule or design vector it proposes. The
feasibility rate is the share of these decisions that satisfy all
constraints when checked in the full model (for example generation
caps, line flows or voltage bounds).
\begin{equation}
  \phi_{\mathrm{feas}}
    = \frac{1}{N}\sum_{j=1}^{N}
      \mathbb{I}\!\big[\mathbf{x}_j \in \mathcal{F}\big],
  \label{eq:feas-rate}
\end{equation}
where $\mathcal{F}$ is the set of decisions that satisfy every
constraint of the full optimization model. When a decision is
infeasible, it also matters how far it exceeds the limits. The largest
breach across the $N$ scenarios and the size of typical large breaches
(for example at the 95th percentile) indicate whether errors remain
small or become critical under high-demand periods and contingencies.
When a constraint must hold only with a stated probability (chance
constraints), the same check is applied on hold-out scenarios: whether
violations occur at most at the promised rate, not only whether the
surrogate error is small on average. Whether a feasible plan is also
near-optimal is assessed separately with the cost-gap metrics in
Section~\ref{sec:val:decision}. Because feasibility reporting is
widespread, only a few representative studies are summarized here,
chosen to differ in domain or constraint type rather than in surrogate
class alone. For unit commitment, learned warm starts are evaluated for
constraint satisfaction on line losses and unit limits
\cite{Chen2022}. In hybrid AC/DC microgrid energy management, hard
constraints are enforced during training of ReLU neural surrogates for
converter losses, and feasibility of the resulting stochastic unit-
commitment solutions is assessed \cite{Liu20242534}. For joint
chance-constrained OPF with renewable generation, feasibility of
quantile-regression surrogates is checked against chance limits
\cite{Chen2023657}. For photovoltaic hosting-capacity planning in
distribution networks, violations of voltage and thermal limits are
reported relative to hosting limits \cite{Koirala20242284}.

\subsection{Decision-aware metrics: regret, cost gap, MIP-gap distribution}
\label{sec:val:decision}

When the surrogate sits inside the optimizer, the validation question
shifts from prediction accuracy (Section~\ref{sec:val:point}) and
constraint satisfaction (Section~\ref{sec:val:feasibility}) to
objective quality: how much extra cost or lost revenue the
surrogate-based plan incurs when evaluated in the full model. As in
Section~\ref{sec:val:feasibility}, each \emph{test scenario} is one
operating condition excluded from surrogate training. On that scenario,
the surrogate-based optimizer returns an operating decision
$\hat{\mathbf{x}}$ and the full model yields a reference optimum
$\mathbf{x}^{*}$. The most common scalar summary is the \emph{relative
cost gap} (also called regret or optimality gap),
Eq.~\eqref{eq:regret}:
\begin{equation}
  r \;=\; \frac{f(\hat{\mathbf{x}}) - f(\mathbf{x}^{*})}{f(\mathbf{x}^{*})},
  \label{eq:regret}
\end{equation}
where $f(\cdot)$ is the objective of the full optimization problem
(for example operating cost or total revenue). Solver-time reduction is
usually reported alongside $r$ because a speed-up without a bounded gap
is not sufficient for acceptance. For mixed-integer problems, studies
also report the \emph{MIP gap distribution} across test scenarios under
a common time limit: the share of instances solved to optimality and
how large the remaining gaps are in the slow tail.

Because cost-gap reporting is now common in learning-to-optimize
studies, only a few representative cases are summarized here; they
do not repeat the feasibility examples above. For large-scale economic
dispatch, optimality gaps and run-time reductions are reported for
end-to-end neural proxies with repair layers \cite{Chen20244723}. For
security-constrained unit commitment with wind and battery storage,
optimality gaps are reported under a fixed solver budget
\cite{Wu2022231}. For joint chance-constrained unit commitment, a
\emph{regret percentage} quantifies the extra cost of the accelerated
method relative to a reference solver \cite{Liang20246508}. In
integrated energy system co-design with a machine-learning market
surrogate, annual revenue error relative to a validation simulation
is reported for price-taker and market-aware formulations
\cite{Jalving2023}. For unit commitment, near-optimality of overall
solutions and average gaps relative to branch-and-cut benchmarks are
reported when machine learning accelerates subproblem solving inside
surrogate Lagrangian relaxation \cite{Wu2024391}. Recent reviews note
that decision-aware metrics are still reported inconsistently and
recommend cost-gap-style summaries as a default
\cite{Ruan2021221,Khaloie2025,Lim2025}.

\subsection{Stress tests: extreme weather, peak load, low-renewable spells}
\label{sec:val:stress}

Decision-aware metrics from Sections~\ref{sec:val:feasibility}
and~\ref{sec:val:decision}, when averaged over a uniform test set,
under-weight the operating regimes that drive cost spikes and
constraint violations. A \emph{stress-test block} is therefore a
deliberate subset of held-out test scenarios---for example
high-demand hours, low-renewable spells, cold or heat waves,
contingencies, or extreme weather that forces line de-rating---on
which the same feasibility, cost-gap or forecast metrics are reported
separately rather than folded into the overall mean.

Because dedicated stress benchmarks are still uncommon, only a few
representative cases are summarized here. Hu et al.\ combine a
Transformer network with Gaussian-process residual learning and report
better net-load forecasts at peaks and during large PV swings on a
hold-out period \cite{Hu2024}. Sun et al.\ optimize Laplace predictive
distributions with pinball loss across 99 wind-power quantiles and
evaluate nine prediction-interval coverage levels on NREL test sites
\cite{Sun20191497}. Du et al.\ train an LSTM surrogate for an
integrated district-heating system with phase-change thermal storage
and report peak heating demand---defined as the average of the three
highest hourly loads---together with price-responsive heating costs
under a peak-shaving strategy \cite{Du2025}. Yang et al.\ use
Gaussian-process renewable intervals to build robust uncertainty sets
and test day-ahead to intra-day microgrid scheduling under worst-case
renewable scenarios \cite{Yang2024__2}. Bayani et al.\ train a
machine-learning surrogate on physics-based line-vibration simulations
and embed it in a wildfire-risk-aware operation model that ranks line
segments for de-energization under extreme wind
\cite{Bayani20231024}. Hu et al.\ emulate stochastic economic
dispatch with a Karhunen--Lo\`eve-reduced Gaussian process and
compare dispatch costs across sampled wind scenarios against Monte
Carlo on the IEEE 118-bus system \cite{Hu2020}. Chen et al.\ train
deep quantile-regression surrogates for joint chance-constrained OPF
and evaluate constraint-violation quantiles under renewable
uncertainty \cite{Chen2023657}. Liang et al.\ report violation rates
when sample-based chance constraints in accelerated unit commitment
are stressed by many extreme renewable scenarios
\cite{Liang20246508}. Khaloie et al., Lim et al.\ and Starke et al.\
consistently note the absence of standardised stress-test designs for
surrogate-assisted energy-system optimization
\cite{Khaloie2025,Lim2025,Starke2025214}.

\subsection{Out-of-distribution and policy-shift checks}
\label{sec:val:ood}

A random hold-out split from the same historical regime tests
interpolation within the training distribution, not extrapolation to
the planning question. An
\emph{out-of-distribution check} therefore evaluates the surrogate on
inputs deliberately drawn from a different regime than the training
data---for example a future year, climate zone, network topology, or
technology portfolio absent during training. A \emph{policy-shift
block} is the planning counterpart: the test scenarios embed a
counterfactual policy or investment decision (new storage, heat pumps,
electrolysers, tariff rules, or market formulations) rather than only
resampling past weather.

Because explicit OOD reporting is still uneven, only a few
representative cases are summarized here. Park et al.\ propose a
physics-induced graph neural network for wind-farm power estimation,
report transferable performance across turbine layouts and wind
conditions, and test it in a layout-optimization experiment
\cite{Park2019}. Subedi et al.\ integrate CNN-based inverter
surrogates into the open-source solver GridLAB-D, validate them
dynamically against reference models, and assess cross-hardware
generalization with transfer learning between inverter platforms
\cite{Subedi2026157}. Volodina et al.\ link Gaussian-process
emulators across a network of municipal energy models and propagate
code uncertainty to unseen input combinations in an OSeMOSYS-based
case study \cite{Volodina2022}. Lawson et al.\ use statistical
emulators in a Bayesian power-network planning framework to quantify
model-output uncertainty at parameter settings that were not included
in the original design-of-experiments \cite{Lawson201647}. Liu
et al.\ design GP kernels for data-driven probabilistic load flow and
select hyperparameters with a potential-game formulation so that the
surrogate generalizes to datasets not used in training
\cite{Liu20222679}. Jalving et al.\ compare price-taker and
market-aware integrated energy system designs optimized with machine-
learning market surrogates against a full validation simulation in
IDAES and report much lower annual revenue error for the market-aware
formulation \cite{Jalving2023}. Khaloie et al., Lim et al., Starke
et al.\ and Mylonopoulos et al.\ recommend that OOD and policy-shift
evidence should accompany surrogate-assisted planning studies
\cite{Khaloie2025,Lim2025,Starke2025214,Mylonopoulos202332697}.

\subsection{Reproducibility: seeds, splits, code, data}
\label{sec:val:repro}

Reported speed-ups, regret values and feasibility rates are only
comparable when the full pipeline can be rerun. Reproducibility in
surrogate-assisted optimization therefore requires, at minimum: the
random seeds used to generate design points and training batches; an
explicit train/validation/test split rule; the surrogate training code;
the host optimization model and solver settings; and access to the
underlying data or a documented procedure to regenerate it. Missing
any one element makes a decision-aware metric hard to interpret or
compare across studies.

Representative positive practice includes generator-parameter
calibration with a stochastic radial-basis-function surrogate and
documented calibration objectives \cite{Khazeiynasab2021}, and fuel-cell
system calibration with an adaptive response surface, Latin hypercube
sampling and stated calibration targets \cite{Zhang2025__7}. Recent
reviews of learning-assisted power-system optimization note uneven
release of code and data, incomplete reporting of splits and seeds,
and inconsistent use of decision-aware metrics, which limits cross-
study comparison
\cite{Ruan2021221,Khaloie2025,Lim2025,Mylonopoulos202332697}. One
review explicitly lists reproducibility, open-source code and
benchmark comparisons among desirable reporting practices for
learning-based OPF \cite{Khaloie2025}. Reproducibility is treated here
as a prerequisite for the metrics in Section~\ref{sec:val:decision}:
a cost gap that cannot be reproduced cannot serve as an acceptance
criterion.
